\section{The degree-seven geometric detector}
\label{sec:geometric-detector}

We give the tensor calculation used in the geometric normalizer
argument, and identify the published representation-theoretic input.
Let $V_g$ be the rational symplectic homology of an oriented genus-$g$
surface with one boundary component. Write $\mathfrak h_{g,1}(k)$
for the degree-$k$ derivations of the free Lie algebra on $V_g$
which annihilate its symplectic boundary element, and let
$J_{g,k}$ be the rational image of the degree-$k$ Johnson homomorphism.
The symplectic identification of a derivation $D$ with its tensor is
\[
 T_D=\sum_{i=1}^g\bigl(x_i\otimes D(y_i)-y_i\otimes D(x_i)\bigr).
\]
Thus $T_D\in V_g^{\otimes(k+2)}$.
For a matching $C$ of positions $1,\ldots,8$, define
\[
 K_C(u_1\otimes\cdots\otimes u_9)
  =\left(\prod_{(i,j)\in C}\langle u_i,u_j\rangle\right)u_9.
\]
Each pairing is ordered with $i<j$.
Morita--Sakasai--Suzuki's five-term detector
\cite[Proposition~4.1, p.~11]{MSS2026} is
\begin{align*}
 R={}&-2K_{(12)(35)(47)(68)}
       +2K_{(12)(36)(47)(58)}\\
    &+2K_{(12)(36)(48)(57)}
       +2K_{(12)(37)(46)(58)}
       +K_{(13)(26)(47)(58)}.
\end{align*}

\begin{proposition}\label{prop:johnson-seven-detector}
For every $g\geq3$, $R(J_{g,7})=0$. The same identity holds
after scalar extension to every $\mathbb Q_\ell$.
\end{proposition}
\begin{proof}
We reconstruct the vanishing from six explicit brackets and the
older stable multiplicity bound in
\cite[proof of Theorem~7.5(iv), author pp.~29--30]{MSS2015}.
That bound says that the standard-representation isotypical
part of $J_{G,7}$ has multiplicity at most six for sufficiently
large $G$. The source obtains it from multiplicities $10,3$
in $\mathfrak h_{G,1}(7)$ and its boundary ideal, Enomoto--Satoh
ranks $3,2$, and the zero intersection in its
Proposition~3.5(ii), author pp.~16--17.

Take $G$ in that stable range, with $G\geq5$, and put
$e_0=x_1,e_1=y_1,\ldots,e_8=x_5,e_9=y_5$.
For an ordered triple $t=(i,j,k)$ set
\[
 \delta_t(z)=\langle z,e_i\rangle[e_j,e_k]
       +\langle z,e_j\rangle[e_k,e_i]
       +\langle z,e_k\rangle[e_i,e_j].
\]
The identification $\bigwedge^3V_G=J_{G,1}$ places these
vectors in the degree-one Johnson image. Use the following
six sequences of triples:
\[
\begin{array}{c|lllllll}
1&(3,7,5)&(1,2,0)&(9,8,6)&(6,8,5)&(2,4,7)&(9,0,3)&(4,6,7)\\
2&(9,5,0)&(6,4,7)&(6,2,0)&(6,5,8)&(7,3,8)&(1,4,3)&(9,2,7)\\
3&(1,3,4)&(4,2,0)&(5,7,0)&(9,2,7)&(5,6,8)&(5,3,0)&(1,4,6)\\
4&(0,1,4)&(3,5,0)&(9,0,5)&(2,3,9)&(4,6,8)&(8,1,0)&(2,7,1)\\
5&(3,7,0)&(4,9,6)&(3,5,2)&(7,3,8)&(6,7,2)&(0,8,9)&(1,2,6)\\
6&(5,4,0)&(0,8,1)&(5,2,3)&(7,4,3)&(0,4,1)&(6,5,2)&(6,7,9)
\end{array}
\]
For row $r$, define $Z_{r,1}=\delta_{t_1}$ and recursively
$Z_{r,j}=[Z_{r,j-1},\delta_{t_j}]$; put $Z_r=Z_{r,7}$.
The Johnson image is a Lie subalgebra, so $Z_r\in J_{G,7}$.
Every tensor term of $Z_r$ has torus weight $x_1$.

In order, take the six contraction matchings
\[
\begin{array}{lll}
(12)(34)(57)(68),&(12)(34)(58)(67),&(12)(35)(47)(68),\\
(12)(36)(47)(58),&(12)(36)(48)(57),&(12)(37)(46)(58).
\end{array}
\]
Their values on $T_{Z_r}$ are multiples of $x_1$, with coefficient
matrix
\[
 A=\begin{pmatrix}
18&12&-2&15&-7&-10\\
0&3&0&-4&4&0\\
15&-24&3&2&-5&6\\
-17&14&-10&11&-11&-4\\
-3&-10&2&-15&3&16\\
2&-16&-1&2&-3&0
\end{pmatrix},\qquad \det A=-203784.
\]
The five-term evaluations give $R(T_{Z_r})=0$ for all six rows.
All these integer entries are obtained by the following finite
recursion, which also specifies the calculation completely.
Represent a word $w=e_{i_1}\cdots e_{i_m}$ in the free associative
algebra, use $[v,w]=vw-wv$, and extend a derivation by
\[
 D(w)=\sum_{j=1}^m e_{i_1}\cdots e_{i_{j-1}}
                    D(e_{i_j})e_{i_{j+1}}\cdots e_{i_m}.
\]
Apply this identity to $[D,E](e_i)=D(E(e_i))-E(D(e_i))$ at each
of the six bracket steps. Form $T_{Z_r}$ by the displayed tensor
identification and apply the four scalar pairings for each matching.
The nonzero tensor supports in the six rows have respectively
$22896,43902,57276,33768,26082,27990$ words. The recurrence gives
$Z_r(\sum_i[x_i,y_i])=0$, the stated weights, the matrix $A$ and
the six zero $R$-values. The program \texttt{check\_johnson\_detector.py}
in the \href{\workbenchurl}{workbench} executes this recursion over
$\mathbb Z$ and records the full verification.

Let $e_V$ be the symplectic-equivariant projection onto the
standard-representation isotypical summand. It preserves $J_{G,7}$.
The contractions and $R$ have target $V_G$, so they kill every other
irreducible summand. The nonzero determinant therefore makes
$e_VZ_1,\ldots,e_VZ_6$ independent in its weight-$x_1$ space.
The published upper bound makes them a basis of that space.
Their symplectic orbits span the whole standard isotypical part.
Equivariance and the six zero values now show that $R$ kills this
part, and its target kills all remaining isotypical parts.
Thus $R(J_{G,7})=0$.

For any given $g\geq3$, choose such a $G\geq g$ and extend a
boundary mapping class by the identity on the added handles.
Its Johnson derivation extends by zero on those handles.
For the symplectic inclusion $i:V_g\hookrightarrow V_G$, the
literal contraction formulas give
\[
 R_G(i^{\otimes9}T)=iR_g(T).
\]
The left side vanishes on the extended Johnson image and $i$ is
injective, proving the assertion at genus $g$. Every step is rational;
scalar extension proves the statement over $\mathbb Q_\ell$, including
primes dividing the displayed integer determinant.
\end{proof}
